\documentclass[slideColor]{prosper}
\usepackage{amsmath,amsfonts,amsthm}
%\usepackage{epsfig}
%\include{Qcircuit}


\newcommand{\bbC}{{\mathbb{C}}}
\newcommand{\C}{{\mathbb{C}}}
\newcommand{\F}{{\mathbb{F}}}
\newcommand{\N}{{\mathbb{N}}}
\newcommand{\Q}{{\mathbb{Q}}}
\renewcommand{\R}{{\mathbb{R}}}
\newcommand{\Z}{{\mathbb{Z}}}

\newcommand{\cD}{{\mathcal{D}}}
\newcommand{\cN}{{\mathcal{N}}}
\newcommand{\cP}{{\mathcal{P}}}
\newcommand{\cQ}{{\mathcal{Q}}}
\newcommand{\cS}{{\mathcal{S}}}
\newcommand{\cU}{{\mathcal{U}}}
\newcommand{\cV}{{\mathcal{V}}}

\newcommand{\tr}{\mathop{\mathrm{tr}}}
\newcommand{\rank}{\mathop{\mathrm{rank}}}
\newcommand{\poly}{\mathop{\mathrm{poly}}}
\newcommand{\E}{\mathop{\mbox{$\mathbf{E}$}}}
\newcommand{\schur}{\mathop{\mathrm{Schur}}}
\newcommand{\schursmall}[2]{{\mathrm{S}}_{#1,#2}}
\newcommand{\planch}{\mathop{\mathrm{Planch}}}
\newcommand{\planchsmall}[1]{{\mathrm{P}}_{#1}}
\renewcommand{\d}{{\mathrm{d}}}

\newcommand{\HSP}{\textsc{hsp}\xspace}

\renewcommand{\>}{\rangle}
\newcommand{\<}{\langle}
\newcommand{\ket}[1]{|#1\rangle}
\newcommand{\bra}[1]{\langle #1|}
\newcommand{\scalar}[2]{\langle #1|#2\rangle}
\newcommand{\proj}[1]{\left|#1\right\>\!\left\<#1\right|}
\newcommand{\norm}[1]{\|#1\|}
\newcommand{\ot}{\otimes}
\newcommand{\eps}{\epsilon}
\newcommand{\ra}{\rightarrow}

\renewcommand{\setminus}{\mathbin{-}}
\newcommand{\partitionof}{\mathbin{\vdash}}


%%%%%%%%%%%%%%%%%%%%%%%%%%%%%%%%%%%%%%%%%%%%%%%%%%%%%%%%%%%%

\title{Weak Fourier-Schur Sampling,\\ the Hidden Subgroup Problem \&\\ the Quantum Collision Problem}
\author{Pawel M. Wocjan\\[1ex]School of Electrical Engineering and Computer Science\\ University of Central Florida\\Orlando}
\email{wocjan@cs.ucf.edu}

\begin{document}

\maketitle

%%%%%%%%%%%%%%%%%%%%%%%%%%%%%%%%%%%%%%%%%%%%%%%%%%%%%%%%%%%%%%%%%%%%%%%%

\begin{slide}{Hidden Subgroup Problem}

\begin{itemize}
\item let $G$ be a finite group and $S$ some finite set
\item let $f : G \rightarrow S$ be a black-box function
\item we have the promise that $f$ hides a subgroup $H\le G$, that is, $f(g)=f(g')$ iff $gH=g'H$
\item the task is to determine the unknown subgroup $H$ (say, in terms of a generating set) as quickly as possible
\item an algorithm is considered to be efficient if it runs in time poly(log($|G|$)
\end{itemize}
\end{slide}

%%%%%%%%%%%%%%%%%%%%%%%%%%%%%%%%%%%%%%%%%%%%%%%%%%%%%%%%%%%%%%%%%%%%%%%%

\begin{slide}{Motivation - Integer Factorization}
\begin{itemize}
\item integer factorization can be reduced (probabilistically) to determining the order of an element $a$ modulo $n$
\item this can be viewed as a HSP over $G:=\Z$
\item let $S:=\Z_n$ and define $f$ by setting $f(x)=a^x$
\item the HSP is $r\Z$, where $r$ is the order of $a$, that is, the smallest positive integer such that $a^r=1$
\end{itemize}
\end{slide}

%%%%%%%%%%%%%%%%%%%%%%%%%%%%%%%%%%%%%%%%%%%%%%%%%%%%%%%%%%%%%%%%%%%%%%%%

\begin{slide}{Motivation - Graph Auto/Isomorphism}
\begin{itemize}
\item graph automorphism (and also graph isomorphism) can be reduced to HSP over the symmetric group $S_n$
\item let $G:=S_n$, $S$ be the set of adjacency matrices of graphs on $n$ vertices, and $A$ be some adjacency matrix
\item define $f$ by setting $f(\pi)=P_{\pi} A P_{\pi}^{-1}$, where $P_{\pi}$ is the permutation matrix of size $n\times n$ corresponding to $\pi$
\item the HSP is the automorphism group of the graph defined by $A$
\end{itemize}
\end{slide}

%%%%%%%%%%%%%%%%%%%%%%%%%%%%%%%%%%%%%%%%%%%%%%%%%%%%%%%%%%%%%%%%%%%%%%%%

\begin{slide}{Classical vs. Quantum Algorithms for HSP}
\begin{itemize}
\item classical query complexity $\Theta(|G|)$
\item quantum query complexity $O(\poly(\log(|G|))$
\item quantum time complexity $O(\poly(\log(|G|))$ for
\begin{itemize}
\item abelian groups
\item Heisenberg groups
\item extraspecial groups
\item and some more (good news: the list has been growing steadily)
\end{itemize}
\item big challenges:
\begin{itemize}
\item symmetric groups $\Longrightarrow$ graph auto/isomorphism
\item dihedral groups $\Longrightarrow$ shortest lattice vector problem 
\end{itemize}
\end{itemize}
\end{slide}

%%%%%%%%%%%%%%%%%%%%%%%%%%%%%%%%%%%%%%%%%%%%%%%%%%%%%%%%%%%%%%%%%%%%%%%%

\begin{slide}{Standard Approach to HSP}
\begin{itemize}
\item evaluate $f$ in superposition
\[
\frac{1}{\sqrt{|G|}} \sum_{g\in G}|g\>|f(g)\>
\]
\item measure second register; assume $s$ is observed; then we obtain the coset state
\[
|gH\>:=\frac{1}{\sqrt{|H|}} \sum_{h\in H}|gh\>
\]
in the first register, where $g$ is such that $f(g)=s$; the element $g$ is completely at random
\end{itemize}
\end{slide}

%%%%%%%%%%%%%%%%%%%%%%%%%%%%%%%%%%%%%%%%%%%%%%%%%%%%%%%%%%%%%%%%%%%%%%%%

\begin{slide}{HSP as Quantum State Identification}
\begin{itemize}
\item using mixed states, this is described by
\[
\rho_H = \frac{1}{|G|}\sum_{g\in G}|gH\>\<gH| 
\]
\item the HSP consists in distinguishing the states $\rho_H$ for the possible $H\le G$
\end{itemize}
\end{slide}

%%%%%%%%%%%%%%%%%%%%%%%%%%%%%%%%%%%%%%%%%%%%%%%%%%%%%%%%%%%%%%%%%%%%%%%%

\begin{slide}{Symmetry of Coset States}
\begin{itemize}
\item the coset state $\rho_H$ can be expressed as
\[
\rho_H = \frac{1}{|G|} \sum_{g\in G} L(g) |H\>\<H| L(g)^\dagger
\]
where $L(g)|h\> = |gh\>$ is the left regular representation of $G$
\item this symmetry can be exploited with the help of Fourier decomposition
\end{itemize}
\end{slide}

%%%%%%%%%%%%%%%%%%%%%%%%%%%%%%%%%%%%%%%%%%%%%%%%%%%%%%%%%%%%%%%%%%%%%%%%

\begin{slide}{Fourier Decomposition}
\begin{itemize}
\item the group algebra $\C G$ decomposes as
\[
  \C G \stackrel{G \times G}{\cong}
  \bigoplus_{\sigma \in \hat G} \cV_\sigma \otimes \cV_\sigma^*
\]
where $\hat G$ denotes a complete set of irreducible representations of $G$, and $\cV_\sigma$ and $\cV_\sigma^*$ are the row and column subspaces acted upon by $\sigma$
\end{itemize}
\end{slide}

%%%%%%%%%%%%%%%%%%%%%%%%%%%%%%%%%%%%%%%%%%%%%%%%%%%%%%%%%%%%%%%%%%%%%%%%

\begin{slide}{Block Structure in the Fourier Basis}
\begin{itemize}
\item $\rho_H$ is invariant under the left multiplication of $G$
\item the Fourier decomposition shows that
\[
\rho_H \cong 
\frac{1}{|G|} 
\bigoplus_{\sigma \in \hat{G}} I_{\dim \cV_\sigma} \otimes \rho_{H,\sigma}
\]
\item this means that $\rho_H$ is block diagonal in the Fourier basis:
\begin{itemize}
\item with blocks labeled by the irreps $\sigma\in\hat G$
\item for each $\sigma$, there is a $\dim \cV_\sigma \times \dim \cV_\sigma$ block $\rho_{H,\sigma}$ that appears $\dim \cV_\sigma$ times (or in word, that it is maximally mixed in the row space)
\end{itemize}
\end{itemize}
\end{slide}

%%%%%%%%%%%%%%%%%%%%%%%%%%%%%%%%%%%%%%%%%%%%%%%%%%%%%%%%%%%%%%%%%%%%%%%%

\begin{slide}{Weak Fourier Sampling}
\begin{itemize}
\item information gain vs. disturbance: measurements extract information about the quantum state, but at the same disturb/destroy it
\item without loss of information, we can 
\begin{itemize}
\item measure the irrep name $\sigma$ and
\item discard the information about which $\sigma$-isotopic block occurred
\end{itemize}
\item the process of measuring the irrep name $\sigma$ is referred to as weak Fourier sampling
\item weak Fourier sampling alone produces insufficient information about $H$ for most nonabelian groups
\end{itemize}
\end{slide}

%%%%%%%%%%%%%%%%%%%%%%%%%%%%%%%%%%%%%%%%%%%%%%%%%%%%%%%%%%%%%%%%%%%%%%%%

\begin{slide}{Strong Fourier Sampling}
\begin{itemize}
\item therefore, a refined measurement must be performed inside the resulting subspace
\item this is referred to as strong Fourier sampling
\item many possibilities; especially if the irrep has large dimension
\end{itemize}
\end{slide}

%%%%%%%%%%%%%%%%%%%%%%%%%%%%%%%%%%%%%%%%%%%%%%%%%%%%%%%%%%%%%%%%%%%%%%%%

\begin{slide}{$k$ copies $\rho_H^{\otimes k}$}
\begin{itemize}
\item just one $\rho_H$ is not sufficient to determine $H$
\item $\Rightarrow$ we must repeat the sampling procedure to obtain statistics
\item however, repeating strong Fourier sampling a polynomial number of times is not sufficient
\item $\Rightarrow$ to solve the HSP in general, we must perform a joint measurement on $k=\poly(\log(|G|))$ copies of $\rho_H^{\otimes k}$
\item in fact, for some groups such as the symmetric group must be entangled across $\Omega(\log(|G|)$ copies
\item $\Rightarrow$ the difficulty of the general HSP may be
attributed at least in part to that fact that highly entangled
measurements are required
\end{itemize}
\end{slide}

%%%%%%%%%%%%%%%%%%%%%%%%%%%%%%%%%%%%%%%%%%%%%%%%%%%%%%%%%%%%%%%%%%%%%%%%

\begin{slide}{Motivation for Schur sampling}

there is another measurement that can also be performed without loss of information

\[
\rho_H \otimes \rho_H \otimes \cdots \otimes \rho_H
\]
we consider the permutation symmetry, that is, that the state $\rho_H^{\otimes k}$ is invariant under permuting the tensor components

\end{slide}

%%%%%%%%%%%%%%%%%%%%%%%%%%%%%%%%%%%%%%%%%%%%%%%%%%%%%%%%%%%%%%%%%%%%%%%%

\begin{slide}{Schur Duality}
\begin{itemize}
\item
the decomposition of $(\C G)^{\otimes k}$ afforded by {\em Schur duality} decomposes $k$ copies of a $d$-dimensional space as 
\[
  (\C^d)^{\otimes k} \stackrel{\cS_k \times \cU_d}{\cong}
  \bigoplus_{\lambda \partitionof k}
  \cP_\lambda \otimes \cQ_\lambda^d
\]
\item the symmetric group $\cS_k$ acts to permute the $k$ registers 
\item the unitary group $\cU_d$ acts identically on each register
\item the subspaces $\cP_\lambda$ and $\cQ_\lambda^d$ correspond to irreps of $\cS_k$ and $\cU_d$, respectively
\item the irreps are labeled by partitions $\lambda \partitionof k$), that is, $\lambda=(\lambda_1,\lambda_2,\ldots)$ where $\lambda_1 \ge \lambda_2 \ge \ldots$ and $\sum_j \lambda_j = k$
\end{itemize}
\end{slide}

%%%%%%%%%%%%%%%%%%%%%%%%%%%%%%%%%%%%%%%%%%%%%%%%%%%%%%%%%%%%%%%%%%%%%%%%

\begin{slide}{Form in Schur basis}
\begin{itemize}
\item $\rho_H^{\otimes k}$ is invariant under the action of $\cS_k$ $\Rightarrow$ the Schur decomposition shows that it is block diagonal 
\item for each $\lambda$, there is a $\dim \cQ_\lambda^{|G|} \times \dim \cQ_\lambda^{|G|}$ block that appears $\dim \cP_\lambda$ times; or  in other words, the state is maximally mixed in the permutation space
\item no information is lost if we measure the partition $\lambda$ and discard the permutation register
\item by analogy to weak Fourier sampling, we refer to this as {\em weak Schur sampling}.  
\item this is a natural measurement to consider (no loss of information and entangling measurement)
\end{itemize}
\end{slide}

%%%%%%%%%%%%%%%%%%%%%%%%%%%%%%%%%%%%%%%%%%%%%%%%%%%%%%%%%%%%%%%%%%%%%%%%


\begin{slide}{Weak Schur Sampling}
\begin{itemize}
\item the distribution under weak Schur sampling is given by
\[
  \Pr(\lambda|\gamma) =
  \tr(\Pi_\lambda \gamma)
\]
\begin{itemize}
\item $\Pi_\lambda$ is the projector onto the $\lambda$-subspace 
\[
  \Pi_\lambda := \frac{\dim \cP_\lambda}{k!} \sum_{\pi \in \cS_k}
                 \chi_\lambda(\pi) \, P(\pi)
\]
\item $\chi_\lambda$ is the character of the irrep of $\cS_k$ labeled
by $\lambda$, and 
\item $P$ is the (reducible) representation of $\cS_k$ that acts to permute the $k$ registers:
$
  P(\pi)|i_1\> \ldots |i_k\> =
  |i_{\pi^{-1}(1)}\>\ldots 
  |i_{\pi^{-1}(k)}\>
$
\end{itemize}

\end{itemize}
\end{slide}

%%%%%%%%%%%%%%%%%%%%%%%%%%%%%%%%%%%%%%%%%%%%%%%%%%%%%%%%%%%%%%%%%%%%%%%%

\begin{slide}{Invariance of the Schur Distribution}
\begin{itemize}
\item the distribution of $\lambda$ according to weak Schur sampling is
invariant under the actions of the permutation and unitary groups:
\[
\Pr(\lambda|\gamma) = 
\Pr(\lambda|P(\pi)U^{\otimes k}\, \gamma\, U^{\dagger \, \otimes k} P(\pi)^\dagger)
\]
for all $U\in\cU_d$ and all $\pi\in\cS_k$
\item in particular, the invariance under $U^{\otimes k}$ implies that for
$\gamma = \rho_H^{\otimes k}$, the distribution according to weak Schur
sampling depends only on the spectrum of $\rho$
\end{itemize}
\end{slide}

%%%%%%%%%%%%%%%%%%%%%%%%%%%%%%%%%%%%%%%%%%%%%%%%%%%%%%%%%%%%%%%%%%%%%%%%

\begin{slide}{Failure of Weak Schur Sampling}
\begin{itemize}
\item the state $\rho_H$ is proportional to a projector of rank $|G|/|H|$
\item suppose we could distinguish between $\rho_H$ for $H=\{1\}$ and some particular $H$ of order $|H| \ge 2$
\item $\Rightarrow$ we could distinguish between 
\begin{itemize}
\item $k$ copies of the maximally mixed state $I_{|G|}/|G|$
\item $k$ copies of the state $J_{|G|/|H|}/(|G|/|H|)$
\end{itemize}
\item $\Rightarrow$ we could distinguish $1$-to-$1$ functions from $|H|$-to-$1$ functions using $k$ queries of the function
\item $\Rightarrow$ this would violate the quantum lower bound for the $|H|$-collision problem stating that $k=\Omega(\sqrt[3]{|G|/|H|})$ copies are required 
\item however, $O(\sqrt[3]{|G|/|H|})$ copies are not sufficient
\end{itemize}

\end{slide}

%%%%%%%%%%%%%%%%%%%%%%%%%%%%%%%%%%%%%%%%%%%%%%%%%%%%%%%%%%%%%%%%%%%%%%%%

\begin{slide}{Quantum Collision Sampling Problem}

{\bf Theorem:}
Given $\rho^{\otimes k}$, distinguishing between
\begin{itemize}
\item $\rho =
I/d$ and 
\item $\rho^2 = \rho/\frac{d}{r}$, that is, $\rho$ is proportional to a projector of rank $d/r$
\end{itemize} 
is possible with success probability $1 - \exp(-\Theta(kr/d))/2$.  

In particular, constant advantage is possible iff $k=\Omega(d/r)$.

\bigskip
In addition to providing the first results on estimation of the spectrum of a quantum state in the regime where $k \ll d^2$, this gives tight estimates of the effectiveness of weak Schur sampling

\end{slide}

%%%%%%%%%%%%%%%%%%%%%%%%%%%%%%%%%%%%%%%%%%%%%%%%%%%%%%%%%%%%%%%%%%%%%%%%

\begin{slide}{Proof Idea}
\begin{itemize}
\item the proof relies on a very careful analysis of the {\em Schur distribution}, $\schur(k,d)$, with
\[
  \Pr(\lambda) = \frac{\dim \cP_\lambda \dim \cQ_\lambda^d}{d^k} 
               = \frac{(\dim \cP_\lambda)^2}{k!}
                  \prod_{(i,j) \in \lambda} \! \left(1 + \frac{j-i}{d}\right)
\]
\item we make use of known results on the typical shape of partitions under the Plancherel distributions 
\item we derive matching lower and upper bounds on the total variation distance of $\schur(k,d)$ and $\schur(k,d/r)$
\end{itemize}
\end{slide}

%%%%%%%%%%%%%%%%%%%%%%%%%%%%%%%%%%%%%%%%%%%%%%%%%%%%%%%%%%%%%%%%%%%%%%%%

\begin{slide}{Failure of Weak Schur Sampling}

{\bf Corollary:} Applying weak Schur sampling to $\rho_H^{\ot k}$ (where $\rho_H$ is defined in, one can distinguish the case $|H|\geq r$ from the case $H=\{1\}$ with constant advantage iff $k=\Omega(|G|/r)$.

\end{slide}

%%%%%%%%%%%%%%%%%%%%%%%%%%%%%%%%%%%%%%%%%%%%%%%%%%%%%%%%%%%%%%%%%%%%%%%%

\begin{slide}{Weak Fourier-Schur Sampling}
\begin{itemize}
\item it is possible to combine Fourier and Schur sampling
\item carry out weak Fourier sampling $k$ times; since the order in which the irreps are obtained does not carry any information, we just consider the type $\underline{\sigma}:=(\sigma_1,\sigma_2,\ldots,\sigma_k) \in \hat G^k$
\item this weak Fourier-type sampling leads to a permutation-symmetric state $\rho_{H,\underline{\sigma}}$
\item apply weak Schur sampling to $\rho_{H,\underline{\sigma}}$
\end{itemize}
it turns out that the distributions of $\underline{\sigma}$ and $\lambda$ are uncorrelated provided that $\underline{\sigma}$ is multiplicity free; this is extremely unlikely for many groups
\end{slide}

\begin{slide}{Failure of Weak Fourier-Schur Sampling}

{\bf Theorem:} The probability that weak Fourier-Schur sampling applied to $\rho_H^{\ot k}$ provides a result that depends on $|H|$ is at most $k^2 d_{\max}^2 |H| / |G|$, where $d_{\max}$ is the largest dimension of an irrep of $G$.

\bigskip

{\bf Corollary (Weak Fourier-Schur sampling on $\cD_N$ and $\cS_n$):}
\begin{itemize}
\item Weak Fourier-Schur sampling on the dihedral group $\cD_N$ cannot
distinguish the trivial subgroup from a hidden reflection with
constant advantage (i.e., success probability $\frac{1}{2} + \Omega(1)$) unless $k = \Omega(\sqrt{N})$. 
\item weak Fourier-Schur sampling on the symmetric group $\cS_n$ or on the
wreath product $\cS_n\wr\Z_2$ cannot
distinguish the trivial subgroup from an order $2$ subgroup with constant
advantage unless $k=\exp(\Omega(\sqrt{n}))$
\end{itemize}

\end{slide}

\end{document}